\chapter{Evenredigheid, snelheid en gemiddelden}\label{ch:g8:speed}

\omterm{def:g7:prop:table}{Evenredigheid} (\cref{ch:g7:prop}) ontmoet nu de fysieke wereld: \omterm{def:g8:speed:def}{snelheid},
debieten, omzettingen van eenheden --- grootheden die \emph{\omterm{def:g3:division:remainder}{quotiënten}} van twee
andere zijn. Het hoofdstuk sluit af met \omterm{def:g8:speed:average}{gewogen gemiddelden}, waar \omterm{def:g6:propdata:def}{evenredig}
denken een klassieke fout voorkomt.

\section{Gemiddelde snelheid}

\begin{definition}[Gemiddelde snelheid]\label{def:g8:speed:def}
De \emph{gemiddelde snelheid}\index{snelheid} van een tocht is het \omterm{def:g3:division:remainder}{quotiënt}
\[
v = \frac{d}{t}
\qquad
\left(\text{afstand gedeeld door tijdsduur}\right),
\]
in kilometer per uur (km/h), meter per seconde (m/s), \dots\ De formule draait
om: $d = v \times t$ en $t = \frac{d}{v}$.
\end{definition}

\begin{example}\label{ex:g8:speed:basic}
Een trein legt $270$ km af in $1$ h $30$ min. Zet eerst de tijd om:
$1$ h $30$ min $= 1.5$ h. Dan is
\[
v = \frac{270}{1.5} = 180 \text{ km/h}.
\]
Afstand bij $90$ km/h tijdens $2$ h $20$ min $= \frac73$ h:
$d = 90 \times \frac73 = 210$ km. Tijd voor $35$ km bij $14$ km/h:
$t = \frac{35}{14} = 2.5$ h $= 2$ h $30$ min.
\end{example}

\begin{remark}[Minuten zijn geen decimalen]
$1$ h $30$ min is $1.5$ h, maar $1$ h $20$ min is \emph{niet} $1.2$ h: het is
$1 + \frac{20}{60} = \frac43 \approx 1.33$ h. Zet minuten altijd om in een \omterm{def:g6:fractions:def}{breuk}
van een uur ($60$ min $= 1$ h) vóór je deelt.
\end{remark}

\begin{omfigure}
\begin{tikzpicture}
\begin{axis}[omaxis,
  width=8.4cm, height=5.6cm,
  xmin=0, xmax=3.4, ymin=0, ymax=270,
  xlabel={tijd (h)}, ylabel={afstand (km)},
  xtick={1,2,3}, ytick={80,160,240}]
  \addplot[omDef, thick, domain=0:3.2] {80*x};
  \addplot[omThm, thick, domain=0:3.2] {50*x};
  \node[omDef, font=\small, rotate=36] at (axis cs:2.35,215)
    {$80$ km/h};
  \node[omThm, font=\small, rotate=25] at (axis cs:2.6,155)
    {$50$ km/h};
\end{axis}
\end{tikzpicture}

{\small Bij constante \omterm{def:g8:speed:def}{snelheid} is de afstand \omterm{def:g6:propdata:def}{evenredig} met de tijd: een rechte
\omterm{def:g6:lines:objects}{lijn} door de oorsprong, waarvan de steilheid de \omterm{def:g8:speed:def}{snelheid} \emph{is}.}
\end{omfigure}

\begin{method}[Snelheden omzetten]\label{met:g8:speed:convert}
Om km/h in m/s om te zetten: $1$ km $= 1000$ m en $1$ h $= 3600$ s, dus is
\[
v \text{ km/h} = \frac{v \times 1000}{3600} \text{ m/s}
= \frac{v}{3.6} \text{ m/s}.
\]
Om m/s in km/h om te zetten, vermenigvuldig je met $3.6$.
\end{method}

\begin{example}\label{ex:g8:speed:convert}
$90$ km/h $= \frac{90}{3.6} = 25$ m/s. Een sprinter die $100$ m in $10$ s loopt,
gaat $10$ m/s $= 36$ km/h.
\end{example}

\section{Grootheden die quotiënten zijn}

\begin{example}[Debiet, prijs per kilo, verbruik]\label{ex:g8:speed:quotients}
De \omterm{def:g8:speed:def}{snelheid} heeft veel neefjes, die je alle op dezelfde manier behandelt:
\begin{itemize}
  \item een kraan vult $48$ L in $4$ min: debiet $\frac{48}{4} = 12$ L/min, dus
        duurt het vullen van een kuip van $150$ L
        $\frac{150}{12} = 12.5$ min;
  \item $0.6$ kg kaas kost $9$ euro: prijs $\frac{9}{0.6} = 15$ euro per kg;
  \item een auto gebruikt $6.3$ L voor $90$ km: verbruik
        $\frac{6.3}{90} \times 100 = 7$ L per $100$ km.
\end{itemize}
Zoek het \omterm{def:g3:division:remainder}{quotiënt}, en de drieledige formule ($q = \frac ab$, $a = qb$,
$b = \frac aq$) doet de \omterm{def:g3:division:remainder}{rest}.
\end{example}

\section{Gewogen gemiddelden}

\begin{definition}[Gewogen gemiddelde]\label{def:g8:speed:average}
Horen bij waarden $v_1, v_2, \dots$ aantallen (of gewichten)
$n_1, n_2, \dots$, dan is hun \emph{gewogen
gemiddelde}\index{gewogen gemiddelde}
\[
\bar v = \frac{n_1 v_1 + n_2 v_2 + \dots}{n_1 + n_2 + \dots} :
\]
het totaal van de waarden, gedeeld door het totaal van de gewichten.
\end{definition}

\begin{example}\label{ex:g8:speed:average}
Twee groepen maakten een toets: groep 1, $12$ leerlingen, gemiddelde $11$; groep
2, $18$ leerlingen, gemiddelde $14$. Gemiddelde van de hele klas:
\[
\bar v = \frac{12 \times 11 + 18 \times 14}{12 + 18}
= \frac{132 + 252}{30} = \frac{384}{30} = 12.8 .
\]
\emph{Niet} $\frac{11 + 14}{2} = 12.5$: de grotere groep trekt het gemiddelde
naar zich toe --- de gewichten doen mee.
\end{example}

\begin{example}[Gemiddelde snelheid over twee stukken]\label{ex:g8:speed:twolegs}
Een wielrenner rijdt $30$ km bij $30$ km/h, en daarna $30$ km bij $15$ km/h. De
\omterm{def:g8:speed:def}{gemiddelde snelheid} over de hele tocht is \emph{niet}
$\frac{30 + 15}{2} = 22.5$ km/h. Reken met de definitie:
\begin{enumerate}
  \item tijden: $\frac{30}{30} = 1$ h, en daarna $\frac{30}{15} = 2$ h; samen
        $3$ h;
  \item totale afstand $60$ km, dus is $v = \frac{60}{3} = 20$ km/h.
\end{enumerate}
Het langzame stuk duurt langer, en weegt dus zwaarder --- een gemiddelde van
snelheden moet gewogen worden met de \emph{tijd}.
\end{example}

\section{Oefeningen}

\begin{exercise}[$\star$]\label{exo:g8:speed:1}
Bereken de \omterm{def:g8:speed:def}{gemiddelde snelheid}: $150$ km in $2$ h; $27$ km in $45$ min; $100$ m
in $12.5$ s.
\end{exercise}

\begin{exercise}[$\star$]\label{exo:g8:speed:2}
Bereken de afstand: $2$ h bij $85$ km/h; $40$ min bij $90$ km/h. Bereken de tijd:
$315$ km bij $90$ km/h (in h en min).
\end{exercise}

\begin{exercise}[$\star$]\label{exo:g8:speed:3}
Zet om: $72$ km/h in m/s; $5$ m/s in km/h; $108$ km/h in m/s.
\end{exercise}

\begin{exercise}[$\star$]\label{exo:g8:speed:4}
Een kraan levert $15$ L/min. Hoe lang duurt het vullen van een tank van $600$ L?
Hoeveel liter in $2$ h $30$?
\end{exercise}

\begin{exercise}[$\star$]\label{exo:g8:speed:5}
Wat is de betere koop: $1.2$ kg appels voor $3.30$ euro, of $0.8$ kg voor $2.16$
euro? Vergelijk de prijzen per kilogram.
\end{exercise}

\begin{exercise}[$\star$]\label{exo:g8:speed:6}
Een klas van $25$ leerlingen heeft een gemiddelde van $12.4$ op een toets; een
andere klas van $35$ leerlingen heeft een gemiddelde van $10$. Bereken het
gemiddelde van de twee klassen samen.
\end{exercise}

\begin{exercise}[$\star$]\label{exo:g8:speed:7}
Geluid gaat ongeveer $340$ m/s. Je ziet een bliksemflits en hoort $6$ seconden
later de donder. Hoe ver sloeg de bliksem in (op $100$ m nauwkeurig)?
\end{exercise}

\begin{exercise}[$\star\star$]\label{exo:g8:speed:8}
Een wandelaar loopt $2$ h bij $5$ km/h, rust $30$ min, en loopt daarna
$1$ h $30$ bij $4$ km/h. Bereken de totale afstand en de \omterm{def:g8:speed:def}{gemiddelde snelheid}
\emph{met de rust erbij} (totale afstand gedeeld door de totaal verstreken tijd).
\end{exercise}

\begin{exercise}[$\star\star$]\label{exo:g8:speed:9}
Cijfers met wegingsfactoren: een leerling haalde $15$ (wegingsfactor $3$), $9$
(wegingsfactor $2$) en $12$ (wegingsfactor $1$). Bereken het \omterm{def:g8:speed:average}{gewogen gemiddelde}.
Welk gewoon (ongewogen) gemiddelde zou ze hebben, en waarom verschillen de twee?
\end{exercise}

\begin{exercise}[$\star\star$]\label{exo:g8:speed:10}
Een auto rijdt $120$ km bij $80$ km/h en daarna $120$ km bij $120$ km/h.
\begin{enumerate}
  \item Bereken de duur van elk stuk, en daarna de \omterm{def:g8:speed:def}{gemiddelde snelheid} over de
        hele rit.
  \item Leg uit waarom de uitkomst kleiner is dan $100$ km/h, het midden van de
        twee snelheden.
\end{enumerate}
\end{exercise}

\begin{exercise}[$\star\star\star$]\label{exo:g8:speed:11}
Twee dorpen liggen $18$ km van elkaar. Anna vertrekt om $10{:}00$ uit het eerste
en wandelt met $5$ km/h naar het tweede; Boris vertrekt op hetzelfde ogenblik uit
het tweede en wandelt met $4$ km/h naar het eerste. Hoe laat ontmoeten ze elkaar,
en op welke afstand van het dorp van Anna? (Samen dichten ze de kloof met
$5 + 4$ km/h.)
\end{exercise}

\section{Opgave: het harmonisch gemiddelde, of waarom de terugrit het gemiddelde
verpest}

\begin{problem}[{Weekendopgave --- de \omterm{def:g8:speed:def}{gemiddelde snelheid} over gelijke afstanden
is $\frac{2 v_1 v_2}{v_1 + v_2}$, en ze haalt het midden van de snelheden nooit}]
\label{pb:g8:speed:1}
\Cref{ex:g8:speed:twolegs} en \cref{exo:g8:speed:10} eindigden beide met dezelfde
verrassing: heen met de ene \omterm{def:g8:speed:def}{snelheid}, terug met de andere, en de \omterm{def:g8:speed:def}{gemiddelde
snelheid} landt \emph{onder} het midden van de twee. Deze opgave vindt de precieze
formule achter die verrassing --- een nieuw soort gemiddelde, het
\emph{harmonisch gemiddelde}\index{harmonisch gemiddelde} --- bewijst dat de
verrassing een stelling is, en drijft ze op tot haar logische uiterste: een
inhaalmanoeuvre die geen enkele \omterm{def:g8:speed:def}{snelheid} ter wereld voor elkaar krijgt.

\medskip
\noindent\textbf{Deel I --- De formule voor een heen-en-terugrit.}
Een tocht legt een afstand $d$ af met \omterm{def:g8:speed:def}{snelheid} $v_1$, en daarna diezelfde afstand
$d$ nog eens met \omterm{def:g8:speed:def}{snelheid} $v_2$ (allemaal positieve getallen).
\begin{enumerate}
  \item Druk de twee tijdsduren uit, en daarna de totale duur, als \omterm{def:g6:fractions:def}{breuken} waarin
        $d$, $v_1$ en $v_2$ voorkomen.
  \item De \omterm{def:g8:speed:def}{gemiddelde snelheid} van de hele tocht is
        $\bar v = \dfrac{2d}{\ \dfrac{d}{v_1} + \dfrac{d}{v_2}\ }$ --- een
        dubbeldekker (\cref{ex:g8:fractions:complex}). Vereenvoudig die en bewijs:
        \[
        \bar v = \frac{2\, v_1 v_2}{v_1 + v_2} .
        \]
  \item Ga de formule na aan beide hierboven aangehaalde berekeningen:
        $v_1 = 30$, $v_2 = 15$ km/h (\cref{ex:g8:speed:twolegs}), en $v_1 = 80$,
        $v_2 = 120$ km/h (\cref{exo:g8:speed:10}).
  \item De afstand $d$ is uit de formule verdwenen. Wat betekent dat concreet
        voor de wielrenner?
  \item Stel nu dat de tocht in de plaats daarvan met elke \omterm{def:g8:speed:def}{snelheid} dezelfde
        \emph{tijd} $T$ doorbrengt. Toon aan dat de \omterm{def:g8:speed:def}{gemiddelde snelheid} dan
        $\frac{v_1 + v_2}{2}$ is, het gewone midden. In de taal van
        \cref{def:g8:speed:average}: met welke grootheid moet een gemiddelde van
        snelheden altijd gewogen worden, en wat verandert er tussen de twee
        scenario's?
\end{enumerate}

\medskip
\noindent\textbf{Deel II --- Harmonisch tegenover rekenkundig.}
Stel voor twee positieve getallen $v_1$ en $v_2$
\[
H = \frac{2\, v_1 v_2}{v_1 + v_2}
\quad\text{(hun \emph{harmonisch gemiddelde}),}
\qquad
A = \frac{v_1 + v_2}{2}
\quad\text{(hun \emph{rekenkundig gemiddelde}).}
\]
\begin{enumerate}[resume]
  \item Bereken $H$ en $A$ voor de paren $(30, 15)$ en $(80, 120)$. Welk van de
        twee gemiddelden wint elke keer?
  \item Zet $A - H$ op de gemeenschappelijke \omterm{def:g4:fractions:def}{noemer} $2(v_1 + v_2)$, werk de
        \omterm{def:g4:fractions:def}{teller} uit met de merkwaardige \omterm{def:g2:mult:def}{producten} (\cref{pb:g8:equations:1}), en
        bewijs:
        \[
        A - H = \frac{(v_1 - v_2)^2}{2\,(v_1 + v_2)} .
        \]
  \item Leid de stelling af die achter alle verrassingen zit: voor positieve
        snelheden is altijd $H \leq A$, met gelijkheid precies als $v_1 = v_2$.
        Waarom beslecht een kwadraat in de \omterm{def:g4:fractions:def}{teller} de zaak?
  \item Bewijs de elegante gelijkheid $H \times A = v_1 \times v_2$: de twee
        gemiddelden vermenigvuldigen terug naar het oorspronkelijke \omterm{def:g2:mult:def}{product}.
  \item Een auto legt $60$ km af bij $40$ km/h en $60$ km bij $60$ km/h.
        Voorspel de \omterm{def:g8:speed:def}{gemiddelde snelheid} met de formule, en bevestig ze daarna
        langs de lange weg, met tijden en totale afstand.
\end{enumerate}

\medskip
\noindent\textbf{Deel III --- De onmogelijke inhaalmanoeuvre.}
\begin{enumerate}[resume]
  \item Een wielrenster start in een wedstrijd van $30$ km met de hoop gemiddeld
        $30$ km/h te halen. Hoeveel tijd mag de wedstrijd haar ten hoogste
        kosten? Ze rijdt de eerste $15$ km bij $15$ km/h: hoeveel van dat
        tijdbudget blijft er over?
  \item Toon aan dat het gehoopte gemiddelde letterlijk onmogelijk is geworden:
        bereken haar \omterm{def:g8:speed:def}{gemiddelde snelheid} als ze de tweede \omterm{def:g3:division:half}{helft} bij $90$ km/h
        rijdt, en leg uit waarom zelfs een willekeurig enorme \omterm{def:g8:speed:def}{snelheid} haar onder
        de $30$ km/h laat.
  \item Ze verlaagt haar doel tot $20$ km/h. Welke \omterm{def:g8:speed:def}{snelheid} op de tweede $15$ km
        haalt dat? Ga je antwoord na met de formule van het \omterm{pb:g8:speed:1}{harmonisch
        gemiddelde}.
  \item Een koudwaterkraan vult een kuip alleen in $20$ min; een warmwaterkraan
        alleen in $30$ min. Welk deel van de kuip vullen de twee kranen samen per
        minuut, en hoe lang doen ze erover? Vergelijk je antwoord met
        $H(20, 30)$ --- wat valt je op?
  \item Een vliegtuig vliegt dezelfde route heen en terug: $900$ km/h met de wind
        mee, $700$ km/h tegen de wind in. Bereken de \omterm{def:g8:speed:def}{gemiddelde snelheid} van de
        heen-en-terugvlucht, en leg in één zin uit waarom wind, hoe hij ook
        waait, een heen-en-terugvlucht altijd \emph{langer} maakt.
\end{enumerate}
\end{problem}
